% ==================== 第6章 ====================
\chapter{实证结果与分析}
\label{ch:empirical}

\section{数据与任务}

本文在四个数据集和一个混合负载上评估FinPlanRAG。四个数据集分别对应长文档确定问答、前瞻性断言裁决、外部迁移和A股年报回测，混合负载用于检验路由。这样的设计使框架的不同主张分别接受检验，而不是用一个数据集的平均结果替代全部证据。

DocFinQA在FinQA问题的基础上加入完整10-K上下文\citep{reddy2024docfinqa}。本文使用测试集中的300题固定样本，答案使用发布的金标准。FinQA程序常量只作为证据召回率的银标准，不作为阅读器输入。该数据集的挑战在于证据分散在长文档中，表格行与段落需要联合阅读，阅读器必须在大量无关内容中定位分子、分母和期间。

收入指引基准包含2010年至2025年美国SEC 8-K文件中的1,005条季度营收指引断言。每条断言包含指引区间、截止时点、当前季度值和证据项。实际季度营收来自SEC XBRL，XBRL没有独立季度时使用披露文件回退。标签为低于区间、区间内、高于区间。数据按时间划分为529条开发集（2010年至2018年）、203条验证集（2019年至2021年）和273条测试集（2022年至2025年）。测试集包含126条高于、110条区间内和37条低于。主结果使用227条具有XBRL核验或XBRL覆盖标签的测试断言；46条纯文本标签单独用于敏感性检验。该数据集的挑战在于标签来自自动核验而非人工标注，类别不平衡明显，且前瞻性断言的真值在截止时点尚未实现。

FinDVer testmini包含700条带财务报告的断言验证样本\citep{zhao2024findver}，作为外部迁移集。本文不在FinDVer上调参，只把在DocFinQA上设计的义务检索策略迁移过去，以检验证据策略是否依赖特定数据分布。

第四个基准是A股年报历史回测。本文选取40家非金融上市公司，覆盖消费、制造、医药、能源和科技行业，对2023年至2025年三个财年构造470条断言，其中营业收入增速232条、归母净利润增速120条、综合毛利率120条、披露日期118条。对每个财年$Y$，问题在$Y+1$年1月31日提出，此时三季报已经披露、年报尚未披露；结果来自随后披露的年报。该基准检验模糊分支在中国市场上的时点判断、数值桥和披露日期断言。

混合负载由300条DocFinQA问题和273条收入指引断言组成，共573条查询。确定查询的效用为词级F1；前瞻性查询的效用为$1-L/10$，其中$L$为决策损失。错误路由记0效用。Oracle路由器把每条查询送到正确分支；Always-RAG把所有查询送到义务检索分支；Always-Anchor把所有查询送到AnchorEvolve。混合效用把两类任务的分数放在同一尺度上，使路由策略可以与固定分支策略比较。

\section{基座模型、具名基线与评价指标}

本文使用三套基座：DeepSeek Flash、GLM 5.3 Flash和Qwen3-8B\citep{qwen2025qwen3}。Qwen3-8B在本地部署。所有方法使用同一基座和相同的提示模板，以比较证据策略和输出契约的差异。

义务检索分支与固定BM25预算1、2、4、8、16段、Self-RAG、CRAG、FLARE、Adaptive-RAG以及义务检索规划器比较。四个具名RAG方法根据原论文描述和公开接口进行同基座适配，保留各自的检索或反思策略，使用与义务检索分支相同的段落块、阅读器和评分，因此差异来自证据策略。

AnchorEvolve与六个具名方法比较：FinDVer \citep{zhao2024findver}、Fin-Fact \citep{rangapur2025finfact}、Self-RAG \citep{asai2024selfrag}、CRAG \citep{yan2024crag}、FLARE \citep{jiang2023flare}和Adaptive-RAG \citep{jeong2024adaptiverag}。每个方法根据原论文描述和公开接口进行同基座适配，通过同一个三分类适配层运行。适配层把原始输出映射为低于、区间内、高于，并提供相同的时点证据块和相同的六条历史断言。学习组件不可用或不兼容时，用使用官方判据的API调用替代，并记录替代关系。六个条件保留各自原论文的判据，不是原系统在其原始数据上的未修改端到端运行。

评价指标分为三组。DocFinQA报告精确匹配、词级F1、答案覆盖率、平均获取段落数、平均输入token和银标准证据召回率。收入指引报告准确率、平衡准确率、宏F1、MCC、AUROC、AUPRC、ECE、Brier、决策损失、可信集覆盖率和平均集合大小。A股回测报告准确率、平衡准确率、宏F1、决策损失、数值MAE、数值MAPE和日期窗口准确率。平衡准确率是各类召回率的平均\citep{brodersen2010balanced,sokolova2009performance}。决策损失使用第\ref{ch:anchor}章参数。损失、ECE和Brier越低越好。置信区间由断言级配对自助法给出，准确率差异使用McNemar检验。路由报告路由准确率、混淆矩阵和混合效用。

评价协议在三个层面控制信息泄漏。第一，收入指引按时间划分，测试年份为2022年至2025年，验证集为2019年至2021年，开发集为2010年至2018年。第二，随机森林使用五个折种子；对每个测试年份$Y$，训练集来自其他公司折，且结果日期早于$Y$年1月1日；测试公司不出现在训练集中。温度缩放在公司不重叠的内部折上拟合。第三，融合权重$w$只在验证集上按决策损失选择，候选为$\{0,0.05,0.1,0.2,0.3,0.5,1.0\}$；最终读出统一使用全动作门，不再在随机森林头与融合头之间切换。测试集不参与拟合、校准、权重选择或动作选择。

所有LLM响应按请求负载的哈希缓存。API密钥只从环境变量读取，不写入结果文件。提示词、运行脚本、逐条预测和表格生成脚本随论文发布。随机森林的折种子、校准方式和融合权重选择在正文中给出。公开结果包含逐条预测、路由决策和汇总指标；LLM响应缓存不随论文提交，但相同输入可以用相同提示词复现。

\section{路由与混合负载}

表\ref{tab:router}报告573条混合负载。DeepSeek Flash路由器正确路由570条（0.995）：300条确定问题中297条，273条前瞻性断言全部正确。GLM 5.3 Flash路由器正确路由567条（0.990）：297条确定问题、270条前瞻性断言。Qwen3-8B正确路由271条确定问题和226条前瞻性断言（0.867），其中29条确定问题被送到AnchorEvolve，47条前瞻性断言被送到义务检索，显示8B基座对契约边界的识别较不稳定。

DeepSeek Flash的混合效用为0.403，与oracle相同，高于Always-RAG的0.121和Always-Anchor的0.282。GLM 5.3 Flash的混合效用为0.419，接近oracle的0.422，并恢复99.3\%的oracle效用。Qwen3-8B的混合效用为0.263，低于Always-Anchor的0.298。结果表明，较大基座能够把确定问答与前瞻性裁决送入相应接口，小基座则需要更强的路由监督或外部信号。

路由错误主要出现在带有未来时态的确定问题，例如“未来支付金额”和“预期股利”。这些问题的答案已经在截止时点文件中披露，但表面表达接近前瞻性断言。该现象说明契约识别依赖事实是否已经存在，而不能仅依据时态或关键词。固定分支策略只能对全部查询采用同一接口，路由器则把契约选择显式化，使路由错误与分支错误可以分别诊断。

\begin{table}[H]
\centering
\zihao{5}
\setlength{\tabcolsep}{3pt}
\caption{混合负载路由结果。效用定义见第\ref{ch:empirical}章。}
\label{tab:router}
\begin{tabular}{lrrrrr}
\toprule
基座 & 路由准确率 & 路由器 & 始终RAG & 始终Anchor & Oracle \\
\midrule
DeepSeek Flash & .995 & .403 & .121 & .282 & .403 \\
GLM 5.3 Flash & .990 & .419 & .130 & .292 & .422 \\
Qwen3-8B & .867 & .263 & .018 & .298 & .316 \\
\bottomrule
\end{tabular}
\end{table}

\section{义务检索在DocFinQA上的结果}

表\ref{tab:docfinqa}报告DocFinQA的300题固定样本。最终义务检索分支使用16段预算，并在阅读器不充分或置信度低于0.8时回退到固定16段证据集。在DeepSeek Flash上，该分支取得0.457精确匹配，平均使用10.69段；固定16段阅读器取得0.390，平均16段；固定8段阅读器取得0.367，平均8段。无回退的预算8规划器取得0.417，平均4.56段；预算12取得0.427，平均5.82段；预算16取得0.460，平均6.85段。回退在24\%的问题上使用更大证据集，弥补了规划器单独的准确率缺口。

在GLM 5.3 Flash上，该分支取得0.497，平均10.31段；固定16段阅读器取得0.480，平均16段；固定8段阅读器取得0.440，平均8段。无回退的预算8规划器取得0.450，平均4.15段；预算12取得0.473，平均5.11段；预算16取得0.470，平均5.88段。回退分支比固定16段阅读器高0.017，平均少用5.69段。两个基座上的答案覆盖率分别为0.770和0.810，高于固定预算的0.640和0.757，说明分支更少弃权。

义务检索分支在准确率和证据成本上优于固定16段基线：准确率更高、证据更少。规划器单独运行时成本更低，但在GLM 5.3 Flash上留下小缺口；置信度门控回退在不对每个问题支付16段成本的情况下补上了缺口。这一结果的含义是证据预算应当由证据状态决定，而不是由查询前固定的$k$决定。对于简单问题，规划器在少量证据后停止；对于复杂问题，回退提供额外证据。

表\ref{tab:docfinqa}同时报告四个具名RAG适配基线。DeepSeek Flash上，Self-RAG与CRAG取得0.373，Adaptive-RAG取得0.363，FLARE取得0.330；GLM 5.3 Flash上，Self-RAG与CRAG取得0.447，Adaptive-RAG与FLARE取得0.393。义务检索分支在DeepSeek Flash上比最好的具名基线高0.084，在GLM 5.3 Flash上高0.050，且平均段落数不高于这些基线。四个条件保留各自的检索或反思策略，使用相同的段落块、阅读器和评分，因此差异来自证据策略而不是阅读器。

Qwen3-8B的DocFinQA结果放在附录\ref{app:qwen}：所有条件的精确匹配都低于0.11，主要原因是阅读器返回解释性文字而不是短答案，该基座不用于义务检索分支的主比较。该结果反映短答案抽取任务对阅读器格式遵循能力的要求，Qwen3-8B的阅读器更偏向解释性输出。框架的分支设计不依赖特定基座，但阅读器的格式遵循能力是实际部署的前提。

\begin{table}[htbp]
\centering
\zihao{5}
\setlength{\tabcolsep}{2pt}
\caption{DocFinQA长文档问答结果。Acc.为精确匹配，Macro F1为宏F1，Paras为平均段落，Cov.为覆盖率。}
\label{tab:docfinqa}
\begin{tabular}{lrrrrrrrr}
\toprule
 & \multicolumn{4}{c}{DeepSeek Flash} & \multicolumn{4}{c}{GLM 5.3 Flash} \\
\cmidrule(lr){2-5}\cmidrule(lr){6-9}
Condition & Acc. & Macro F1 & Paras & Cov. & Acc. & Macro F1 & Paras & Cov. \\
\midrule
Fixed-4 & .333 & .192 & 4.00 & .557 & .400 & .199 & 4.00 & .653 \\
Fixed-8 & .367 & .223 & 8.00 & .603 & .440 & .225 & 8.00 & .720 \\
Fixed-16 & .390 & .234 & 16.00 & .640 & .480 & .249 & 16.00 & .757 \\
Planner B=8 & .417 & .230 & 4.56 & .707 & .450 & .248 & 4.15 & .763 \\
Planner B=16 & .460 & .276 & 6.85 & .773 & .470 & .243 & 5.88 & .840 \\
Self-RAG & .373 & .223 & 8.03 & .613 & .447 & .231 & 8.16 & .717 \\
CRAG & .373 & .222 & 8.05 & .613 & .447 & .235 & 8.32 & .727 \\
FLARE & .330 & .197 & 4.00 & .557 & .393 & .193 & 4.19 & .643 \\
Adaptive-RAG & .363 & .211 & 7.22 & .570 & .393 & .198 & 6.51 & .657 \\
\midrule
\textbf{FinPlanRAG (obligation + fallback)} & \textbf{.457} & \textbf{.276} & 10.69 & \textbf{.770} & \textbf{.497} & \textbf{.252} & 10.31 & \textbf{.810} \\
\bottomrule
\end{tabular}
\end{table}


总体来看，义务检索分支在DeepSeek Flash和GLM 5.3 Flash上均以低于固定16段的平均证据量取得更高的准确率和覆盖率。规划器单独运行的成本最低，回退主要改善规划器无法闭合的难题；两者组合后的分支在准确率、覆盖率和证据成本之间取得较优的平衡。与四个具名RAG适配基线相比，义务检索的优势来自证据策略，而不是来自不同的阅读器或评分口径。

\section{前瞻性断言裁决结果}

表\ref{tab:main}按方法列出三套基座的结果。纵向是方法，横向是基座和指标，每个基座给出准确率、平衡准确率、宏F1和全动作决策损失。DeepSeek使用验证集接受的扩展类型库，GLM和Qwen保留基础类型库。

\begin{table}[htbp]
\centering
\zihao{-5}
\setlength{\tabcolsep}{2.2pt}
\caption{主结果（227条核验断言）。Acc.为准确率，Bal.为平衡准确率，Macro为宏F1，Loss为全动作决策损失。}
\label{tab:main}
\begin{tabular}{lrrrrrrrrrrrrr}
\toprule
\multirow{2}{*}{方法} & \multicolumn{4}{c}{DeepSeek Flash} & \multicolumn{4}{c}{GLM 5.3 Flash} & \multicolumn{4}{c}{Qwen3-8B} \\
\cmidrule(lr){2-5}\cmidrule(lr){6-9}\cmidrule(lr){10-13}
 & Acc. & Bal. & Macro & Loss & Acc. & Bal. & Macro & Loss & Acc. & Bal. & Macro & Loss \\
\midrule
FinDVer & .515 & .424 & .423 & 3.128 & .507 & .456 & .458 & 3.198 & .463 & .441 & .429 & 5.339 \\
Fin-Fact & .524 & .428 & .429 & 3.251 & .489 & .401 & .400 & 3.410 & .507 & .473 & .459 & 4.881 \\
Self-RAG & .546 & .460 & .468 & 3.031 & .507 & .454 & .460 & 3.181 & .485 & .446 & .431 & 5.119 \\
CRAG & .515 & .412 & .407 & 3.004 & .498 & .424 & .427 & 3.480 & .467 & .433 & .417 & 5.251 \\
FLARE & .511 & .418 & .412 & 3.278 & .489 & .407 & .408 & 3.559 & .326 & .298 & .212 & 6.661 \\
Adaptive-RAG & .489 & .402 & .396 & 3.612 & .436 & .366 & .366 & 4.070 & .352 & .325 & .290 & 6.396 \\
\midrule
LLM-only & .498 & .397 & .362 & 3.445 & .533 & .428 & .414 & 3.427 & .410 & .357 & .267 & 5.498 \\
RF & .533 & .491 & .486 & 2.987 & .533 & .491 & .486 & 2.987 & .533 & .491 & .486 & 2.987 \\
\textbf{FinPlanRAG (Ours)} & \textbf{.577} & \textbf{.526} & \textbf{.530} & \textbf{2.907} & \textbf{.551} & \textbf{.504} & \textbf{.503} & \textbf{2.943} & \textbf{.542} & \textbf{.501} & \textbf{.499} & \textbf{2.925} \\
\bottomrule
\end{tabular}
\end{table}

表中的FinPlanRAG指锚点后验与随机森林后验按验证集权重融合后的系统，融合权重只在验证集选择。DeepSeek的FinPlanRAG使用验证集扩展后的类型库，GLM和Qwen使用基础类型库；类型库扩展是训练步骤，不是另一个预测模块。FinPlanRAG在三个基座的准确率、平衡准确率、宏F1和决策损失上均为最高。DeepSeek的FinPlanRAG相对最强基线CRAG把损失从3.004降到2.907，准确率从.515提高到.577；GLM相对Self-RAG把损失从3.181降到2.943，准确率从.507提高到.551；Qwen相对Fin-Fact把损失从4.881降到2.925，准确率从.507提高到.542。配对自助法给出的损失差分别为$-0.097$ $[-0.502,0.300]$、$-0.238$ $[-0.661,0.176]$和$-1.956$ $[-2.617,-1.322]$，只有Qwen的区间排除零。

{\sloppy FinPlanRAG相对随机森林的损失差为$-0.079$ $[-0.300,0.141]$、$-0.044$ $[-0.203,0.115]$和$-0.062$ $[-0.282,0.159]$，三个区间均包含零。\par}融合在点估计上低于随机森林，但总体差异不足以主张稳定优势。把测试集按真实结果分层后，融合在区间内样本上的优势更清楚，见第\ref{sec:within}节。

表\ref{tab:calibration}给出概率质量指标。融合系统在三个基座上的AUROC为.713、.711和.707，ECE为.064、.030和.069，Brier为.555、.557和.561；对应的随机森林AUROC为.706，ECE为.033，Brier为.566。FinPlanRAG的排序能力和Brier分数优于随机森林；ECE在DeepSeek和Qwen上略高，在GLM上略低，说明融合后验的排序更准，但置信度校准仍有改进空间。

\begin{table}[htbp]
\centering
\zihao{5}
\setlength{\tabcolsep}{4pt}
\caption{概率质量与校准（227条核验断言）。AUROC和AUPRC越高越好，ECE和Brier越低越好。}
\label{tab:calibration}
\begin{tabular}{lrrrrrrrrr}
\toprule
\multirow{2}{*}{方法} & \multicolumn{3}{c}{DeepSeek Flash} & \multicolumn{3}{c}{GLM 5.3 Flash} & \multicolumn{3}{c}{Qwen3-8B} \\
\cmidrule(lr){2-4}\cmidrule(lr){5-7}\cmidrule(lr){8-10}
 & AUROC & ECE & Brier & AUROC & ECE & Brier & AUROC & ECE & Brier \\
\midrule
LLM-only & .668 & .110 & .612 & .651 & .101 & .609 & .493 & .439 & .975 \\
RF & .706 & .033 & .566 & .706 & .033 & .566 & .706 & .033 & .566 \\
\textbf{FinPlanRAG (Ours)} & \textbf{.713} & \textbf{.064} & \textbf{.555} & \textbf{.711} & \textbf{.030} & \textbf{.557} & \textbf{.707} & \textbf{.069} & \textbf{.561} \\
\bottomrule
\end{tabular}
\end{table}

\section{决策读出的作用}

表\ref{tab:readout}在同一后验上比较受限可信集门与全动作读出。受限门只比较点答案和最高密度可信集，全动作门在点答案、非单元素裁决集和弃权上选择期望损失最小的动作。融合后验在三个基座上的全动作损失相对受限门分别下降0.749、0.581和0.696，95\%配对区间均排除零。随机森林后验的下降为0.467，区间同样排除零。全动作读出没有改变后验，也没有增加训练数据，它改变的是给定后验时的动作选择。

\begin{table}[htbp]
\centering
\zihao{-5}
\setlength{\tabcolsep}{2.5pt}
\caption{受限门与全动作读出的决策损失（227条核验断言）。$\Delta$为全动作减受限门。}
\label{tab:readout}
\begin{tabular}{lrrrrrrrrr}
\toprule
\multirow{2}{*}{方法} & \multicolumn{3}{c}{DeepSeek} & \multicolumn{3}{c}{GLM} & \multicolumn{3}{c}{Qwen3} \\
\cmidrule(lr){2-4}\cmidrule(lr){5-7}\cmidrule(lr){8-10}
 & 受限门 & 全动作 & $\Delta$ & 受限门 & 全动作 & $\Delta$ & 受限门 & 全动作 & $\Delta$ \\
\midrule
\textbf{FinPlanRAG (Ours)} & \textbf{3.656} & \textbf{2.907} & \textbf{$-0.749$} & \textbf{3.524} & \textbf{2.943} & \textbf{$-0.581$} & \textbf{3.621} & \textbf{2.925} & \textbf{$-0.696$} \\
RF & 3.454 & 2.987 & $-0.467$ & 3.454 & 2.987 & $-0.467$ & 3.454 & 2.987 & $-0.467$ \\
LLM-only & 3.674 & 3.445 & $-0.229$ & 3.454 & 3.427 & $-0.026$ & 5.048 & 5.498 & $+0.449$ \\
\bottomrule
\end{tabular}
\end{table}

全动作读出的作用集中在高不确定性区域。融合后验在227条测试断言中有217条选择裁决集、10条选择点答案、0条弃权；平均集合大小为2.12。点答案不是唯一合理的输出，当两个相邻结果的期望损失接近时，返回二元裁决集可以降低误判成本。这一结果说明，前瞻性断言的核心问题不只是预测哪个标签，还包括给定后验应当输出什么。

\section{指导区间内的分层结果}
\label{sec:within}

总体损失打平时，锚点层是否在某个误差层上提供互补信号，是一个独立问题。表\ref{tab:within}报告真实结果落在指引区间内的82条测试断言。FinPlanRAG在三个基座上均低于随机森林：DeepSeek从3.146降到2.780，GLM从3.146降到2.902，Qwen从3.146降到2.732（以第一个随机森林种子为例）。五个随机森林种子中，DeepSeek和Qwen的五个区间全部排除零，GLM有四个区间排除零。

\begin{table}[htbp]
\centering
\zihao{5}
\setlength{\tabcolsep}{4pt}
\caption{指导区间内样本的分层比较（82条测试断言）。差值为FinPlanRAG减随机森林，负值表示FinPlanRAG损失更低。}
\label{tab:within}
\begin{tabular}{lrrrr}
\toprule
基座 & FinPlanRAG & 随机森林 & 损失差 [95\% CI] & 种子一致性 \\
\midrule
DeepSeek & \textbf{2.780} & 3.146 & $-0.366$ $[-0.732,-0.049]$ & 5/5 \\
GLM & \textbf{2.902} & 3.146 & $-0.244$ $[-0.537,-0.024]$ & 4/5 \\
Qwen3 & \textbf{2.732} & 3.146 & $-0.415$ $[-0.732,-0.170]$ & 5/5 \\
\bottomrule
\end{tabular}
\end{table}

这个分层结果说明锚点层在高密度、低极端误差的区域提供了互补信号。分层依据是真实结果，预测时不可知，因此它只能解释误差来源，不能作为部署时的路由规则。总体损失上的差异仍按不显著报告。

\section{A股年报历史回测}
\label{sec:ashare}

表\ref{tab:ashare}给出A股断言的分类与决策结果，DeepSeek和GLM使用470条全样本，Qwen3使用250条分层子样本。上年外推基线把上年增速或上年毛利率直接外推，季度年化基线把三季报年化后外推；完整系统使用锚点、资格判定、时间演化、数值桥和全动作读出，去桥变体关闭数值桥。

\begin{table}[htbp]
\centering
\zihao{-5}
\setlength{\tabcolsep}{3pt}
\caption{A股年报历史回测。Acc.为准确率，Bal.为平衡准确率，Macro为宏F1，Loss为全动作决策损失，日期为披露日期窗口准确率。Qwen3在250条分层子样本上评估，其余方法在470条全样本上评估。}
\label{tab:ashare}
\begin{tabular}{llrrrrr}
\toprule
方法 & 样本 & Acc. & Bal. & Macro & Loss & 日期 \\
\midrule
\multicolumn{7}{l}{\textit{传统基线}} \\
上年外推 & 470 & .491 & .402 & .401 & 5.094 & .780 \\
季度年化 & 470 & .777 & .786 & .785 & 2.270 & .729 \\
\midrule
\multicolumn{7}{l}{\textit{DeepSeek}} \\
\textbf{FinPlanRAG} & 470 & \textbf{.783} & \textbf{.778} & \textbf{.776} & \textbf{2.179} & \textbf{.788} \\
去数值桥 & 470 & .638 & .591 & .590 & 2.421 & .788 \\
\midrule
\multicolumn{7}{l}{\textit{GLM}} \\
\textbf{FinPlanRAG} & 470 & \textbf{.549} & \textbf{.485} & \textbf{.483} & \textbf{3.332} & \textbf{.754} \\
去数值桥 & 470 & .549 & .485 & .483 & 3.332 & .754 \\
\midrule
\multicolumn{7}{l}{\textit{Qwen3}} \\
季度年化 & 250 & .772 & .789 & .789 & 2.324 & .714 \\
\textbf{FinPlanRAG} & 250 & \textbf{.580} & \textbf{.523} & \textbf{.522} & \textbf{4.208} & \textbf{.690} \\
去数值桥 & 250 & .564 & .500 & .499 & 3.864 & .690 \\
\bottomrule
\end{tabular}
\end{table}

DeepSeek基座上的完整系统取得.783的准确率和2.179的决策损失，高于季度年化的.777和2.270，也明显高于去桥变体的.638和2.421。完整系统相对去桥变体的损失差为$-0.243$ $[-0.404,-0.077]$，准确率差为$+0.145$ $[0.109,0.183]$，两个区间均排除零。数值桥把营业收入和归母净利润的增速判断从标签后验推进到可计算的数值区间，是DeepSeek增益的主要来源。GLM基座上完整系统与去桥变体相同，说明该基座的抽取结果没有提供足够的数值锚点，桥接未触发；完整系统也低于季度年化。Qwen3的250条子样本上，完整系统的准确率比去桥变体高$+0.016$ $[0.004,0.032]$，但损失高$+0.344$ $[0.064,0.640]$，说明数值桥改善了部分区间判断，却没有改善整体决策校准。A股结果因此分为两类：DeepSeek上数值桥带来明确增益，GLM和Qwen上抽取质量不足以支撑桥接，标签后验仍是更稳的回退。

表\ref{tab:ashare-numeric}给出数值桥的点估计。DeepSeek完整系统的数值MAE为9.86个百分点，季度年化为7.40个百分点；完整系统的分类损失更低，但数值点估计不如简单年化。这说明两种目标并不等价：数值桥改善的是区间判断和决策损失，不保证点值误差最小。Qwen3完整系统的数值MAE为39.21个百分点，明显高于DeepSeek，说明其抽取器输出的数值字段噪声较大。GLM没有产生可用的数值点估计，数值桥未触发。披露日期窗口准确率见表\ref{tab:ashare}：DeepSeek为.788，GLM为.754，Qwen3为.690；对应的季度年化分别为.729、.729和.714。

\begin{table}[htbp]
\centering
\zihao{5}
\setlength{\tabcolsep}{4pt}
\caption{A股回测的数值桥点估计。MAE为增速或毛利率的百分点误差，MAPE只在真实值非零时计算。GLM没有产生可用的数值点估计，数值桥未触发。}
\label{tab:ashare-numeric}
\begin{tabular}{llrr}
\toprule
方法 & 样本 & 数值MAE & 数值MAPE \\
\midrule
季度年化 & 470 & 7.40 & 459.7 \\
\textbf{FinPlanRAG (DeepSeek)} & 470 & \textbf{9.86} & \textbf{408.7} \\
\textbf{FinPlanRAG (Qwen3)} & 250 & \textbf{39.21} & \textbf{508.7} \\
\bottomrule
\end{tabular}
\end{table}

\section{公开基准的组件级检查}

表\ref{tab:findver-component}给出FinDVer前150条断言的结果。状态条件式获取从两条证据开始，只有在阅读器报告证据不足时才扩展，最多八条；固定预算条件直接给4条或8条；金标条件给出专家标注的证据段落。

\begin{table}[H]
\centering
\zihao{5}
\setlength{\tabcolsep}{4pt}
\caption{FinDVer组件检查（150条断言）。Acc.为断言判定准确率，Macro为宏F1，Recall为证据段落召回率，Para.为平均证据段落数。}
\label{tab:findver-component}
\zihao{-5}
\setlength{\tabcolsep}{2.2pt}
\begin{tabular}{lrrrrrrrrrrrr}
\toprule
\multirow{2}{*}{条件} & \multicolumn{4}{c}{DeepSeek} & \multicolumn{4}{c}{GLM} & \multicolumn{4}{c}{Qwen3} \\
\cmidrule(lr){2-5}\cmidrule(lr){6-9}\cmidrule(lr){10-13}
 & Acc. & Macro & Recall & Para. & Acc. & Macro & Recall & Para. & Acc. & Macro & Recall & Para. \\
\midrule
固定4条 & .767 & .550 & .620 & 4.00 & .760 & .485 & .620 & 4.00 & .733 & .438 & .620 & 4.00 \\
固定8条 & .847 & .570 & .755 & 8.00 & .860 & .573 & .755 & 8.00 & .700 & .426 & .755 & 8.00 \\
金标证据 & .913 & .577 & 1.000 & 2.99 & .933 & .603 & 1.000 & 2.99 & .707 & .441 & 1.000 & 2.99 \\
\textbf{FinPlanRAG 获取} & \textbf{.873} & \textbf{.559} & \textbf{.603} & \textbf{4.63} & \textbf{.913} & \textbf{.592} & \textbf{.628} & \textbf{4.95} & \textbf{.660} & \textbf{.397} & \textbf{.517} & \textbf{4.01} \\
\bottomrule
\end{tabular}
\end{table}

DeepSeek和GLM上，FinPlanRAG的状态条件式获取在约一半段落数下取得高于固定8条的判定准确率：DeepSeek为.873对.847，GLM为.913对.860。证据召回率低于固定8条，说明它没有覆盖全部金标段落，但覆盖了足以判定的段落。Qwen3上的方向相反：固定4条为.733，状态条件式为.660，低于金标条件.707。Qwen3的阅读器在短上下文下更稳定，扩展证据反而引入噪声。这个结果把状态条件式获取的适用条件限定在阅读器能够利用新增证据的基座上。



% ==================== 第7章 消融实验 ====================
\section{消融实验}

本章的消融在锚点分支和最终系统上分别进行。锚点分支的损失只用于比较证据策略，不与第\ref{ch:empirical}章的融合主表直接比较；最终系统的消融使用与主表相同的后验和读出。

\subsection{锚点提示词与时间演化}

文本锚点的作用首先出现在纯LLM分支。旧提示词使用结构化摘要，新提示词使用时点原文、段落编号和逐字引文。DeepSeek上的准确率、平衡准确率和宏F1分别从.445、.307、.242提高到.502、.397、.364。来源落地和逐字引文让模型能够区分“指引区间”“历史实际”和“分析师预测”，减少了把累计增速当作全年增速的错误。

时间演化的作用较小但方向一致。新锚点从静态可靠度切换到按揭示日更新的Beta可靠度，损失从2.908降到2.855；旧锚点从2.899降到2.793。准确率变化不大，演化主要改善概率校准和决策损失，而不是点分类。这个结果符合演化的设计目标：它更新的是锚点类型的可靠性，不是标签本身。

\subsection{资格判定}

表\ref{tab:qual-ablation}比较完整资格判定与接受全部解析锚点。完整判定拒绝空泛、过期、对抗和不合格的锚点，接受率在三个基座上分别为27.6\%、42.7\%和68.3\%；接受全部条件保留所有解析结果。

\begin{table}[htbp]
\centering
\zihao{5}
\setlength{\tabcolsep}{4pt}
\caption{资格判定消融（227条核验断言，锚点分支）。差值为接受全部减完整判定，正值表示完整判定损失更低。}
\label{tab:qual-ablation}
\begin{tabular}{lrrrr}
\toprule
基座 & 完整判定 & 接受全部 & 损失差 [95\% CI] & 接受率 \\
\midrule
DeepSeek & 3.304 & 3.507 & $+0.203$ $[-0.106,0.520]$ & 27.6\% \\
GLM & 3.339 & 3.427 & $+0.088$ $[-0.176,0.361]$ & 42.7\% \\
Qwen3 & 4.599 & 4.652 & $+0.053$ $[-0.141,0.256]$ & 68.3\% \\
\bottomrule
\end{tabular}
\end{table}

三个基座的点估计都支持完整判定，DeepSeek的增益最大，Qwen3的增益最小。三个配对区间都包含零，因此资格判定的作用按方向性证据报告。Qwen3的接受率明显高于另外两个基座，说明其抽取器产生了更多边缘锚点，资格判定拒绝的比例更高，但剩余锚点的区分力有限。

\subsection{融合权重}

表\ref{tab:weight}给出验证集上不同融合权重的决策损失。权重$w$是锚点后验在融合后验中的占比，$w=0$对应纯随机森林，$w=1$对应纯锚点后验。DeepSeek在$w=0.2$取得最低验证损失，GLM和Qwen3在$w=0.1$取得最低验证损失；测试集使用验证集选择的权重，不参与调参。

\begin{table}[htbp]
\centering
\zihao{5}
\setlength{\tabcolsep}{4pt}
\caption{融合权重的验证集损失（185条验证断言）。粗体为验证集选择。}
\label{tab:weight}
\begin{tabular}{lrrrrrrr}
\toprule
基座 & 0.0 & 0.05 & 0.1 & 0.2 & 0.3 & 0.5 & 1.0 \\
\midrule
DeepSeek & 3.341 & 3.308 & 3.265 & \textbf{3.232} & 3.341 & 3.503 & 3.859 \\
GLM & 3.341 & 3.211 & \textbf{3.200} & 3.286 & 3.362 & 3.416 & 3.708 \\
Qwen3 & 3.341 & 3.362 & \textbf{3.330} & 3.686 & 3.859 & 4.162 & 6.519 \\
\bottomrule
\end{tabular}
\end{table}

三个基座的最优权重都偏向随机森林，DeepSeek为0.2，GLM和Qwen3为0.1。锚点后验单独使用时损失明显更高，Qwen3在$w=1$时达到6.519。融合的价值来自补充信号，而不是替代结构化后验。

\subsection{类型库的验证集扩展}

类型库扩展是算法的一个训练步骤，不是独立的预测模块。算法\ref{alg:typeexpand}先在验证集上挖掘高损失断言，由LLM提出候选类型，再在选择集上重抽锚点并计算损失；只有验证损失下降超过$\delta$的候选才进入类型库。接受后的类型库冻结，测试集只评估一次。


DeepSeek上的候选类型库包含“指引中点与当前季度收入之比”和“财年周数效应”两个类型。验证选择集上，基础类型库的损失为3.943，候选类型库为3.792，下降0.151，超过阈值，候选被接受。提示词单独修改和类型加提示词组合两个候选没有通过验证门控。冻结类型库后，测试集上的准确率从.555提高到.577，宏F1从.513提高到.530，损失从2.934降到2.907；配对损失差为$-0.026$ $[-0.159,0.097]$，区间包含零。

表\ref{tab:typeevolve}汇总这一过程。类型库扩展的测试增益是小的，但方向与验证集一致，且没有使用测试标签。GLM和Qwen3没有候选通过验证门控，保留基础类型库，因此主表中只有DeepSeek使用扩展类型库。

\begin{table}[htbp]
\centering
\zihao{5}
\setlength{\tabcolsep}{4pt}
\caption{DeepSeek类型库扩展（227条核验断言）。Loss为全动作决策损失。}
\label{tab:typeevolve}
\begin{tabular}{lrrrr}
\toprule
系统 & Acc. & Bal. & Macro & Loss \\
\midrule
基础类型库 & .555 & .511 & .513 & 2.934 \\
扩展类型库 & .577 & .526 & .530 & 2.907 \\
差值 & $+.022$ & $+.015$ & $+.017$ & $-0.026$ $[-0.159,0.097]$ \\
\bottomrule
\end{tabular}
\end{table}

\subsection{去掉锚点结构的对照}
\label{sec:no-anchor}

前面的消融检验的是组件内部，例如资格判定是否比接受全部更好。一个更根本的问题是：如果去掉三层锚点结构，只保留时点证据、LLM标签分数、温度缩放和全动作读出，损失会差多少。这个对照直接检验锚点结构本身的必要性。

本文从同一批锚点抽取响应中只保留标签分数，去掉全部锚点对象、资格判定和Beta可靠性，在验证集上拟合温度参数，再在测试集上使用同一个全动作读出。这个对照可以看作FinAbstain式选择性预测在本任务上的近似：时点证据、概率输出、温度缩放和损失感知弃权都在，但没有多类别资格判定和类型可靠性演化。表\ref{tab:noanchor}给出结果。去掉锚点结构后，DeepSeek、GLM和Qwen的损失分别上升0.379、0.317和0.899，准确率分别下降5.3、8.8和13.7个百分点。时点证据和选择性输出仍然存在，性能却明显低于完整系统，说明增益不能只归因于“时点检索加弃权”这一一般性组合。

\begin{table}[htbp]
\centering
\zihao{5}
\setlength{\tabcolsep}{4pt}
\caption{去掉三层锚点结构的架构级对照（227条核验断言）。温度在验证集上选择，测试集只评估一次。}
\label{tab:noanchor}
\begin{tabular}{lrrrrrrrrr}
\toprule
\multirow{2}{*}{方法} & \multicolumn{3}{c}{DeepSeek} & \multicolumn{3}{c}{GLM} & \multicolumn{3}{c}{Qwen3} \\
\cmidrule(lr){2-4}\cmidrule(lr){5-7}\cmidrule(lr){8-10}
 & Acc. & Macro & Loss & Acc. & Macro & Loss & Acc. & Macro & Loss \\
\midrule
无锚点结构 & .524 & .398 & 3.286 & .463 & .347 & 3.260 & .405 & .257 & 3.824 \\
\textbf{FinPlanRAG} & \textbf{.577} & \textbf{.530} & \textbf{2.907} & \textbf{.551} & \textbf{.504} & \textbf{2.943} & \textbf{.542} & \textbf{.501} & \textbf{2.925} \\
\bottomrule
\end{tabular}
\end{table}

这个对照也解释了为什么本文把贡献放在证据资格判定和类型可靠性演化上，而不是放在“选择性预测”上。选择性预测本身是一个一般性工具，FinAbstain等工作已经把它用于金融预测；本文的增量在于锚点结构把证据拆成可判定、可追溯、可按类型更新的对象。去掉这层结构，系统仍然可以选择弃权或裁决集，但失去了区分证据失效模式的能力。

\subsection{数值桥}

数值桥的作用在A股回测上最清楚。DeepSeek完整系统的损失为2.179，去桥变体为2.421，差值为$-0.242$；准确率从.638提高到.783，宏F1从.590提高到.776。桥接把“增速可能高于区间”的标签判断转化为“按当前季度年化，全年增速约为多少”的数值计算，减少了纯标签后验的过度平滑。

桥接的收益依赖抽取质量。GLM完整系统与去桥变体完全相同，说明该基座没有产生足够的数值锚点，桥接未触发；Qwen3的结果见第\ref{ch:empirical}章。数值桥不是无条件增益，它要求抽取器输出可核验的数值字段、单位和期间。对抽取质量较弱的基座，标签后验仍是更稳的回退。

\subsection{消融结论}

消融结果可以归纳为四点。第一，去掉三层锚点结构后，三个基座的损失分别上升0.379、0.317和0.899，说明锚点结构本身提供了不可替代的增量，而不是只有选择性输出在起作用。第二，输出形式和资格判定是稳定的组件：全动作读出在三个基座上显著降低损失，资格判定方向一致但区间包含零。第三，类型库扩展和数值桥是条件增益：它们在DeepSeek上有效，在抽取质量较弱的基座上不触发或收益有限。第四，融合权重偏向随机森林，锚点层提供的是补充信号，而不是替代结构化后验。这个结论与主表的定位一致：框架的优势来自证据状态、决策接口和可审计性，而不是单一组件的普遍提升。


\section{外部迁移与误差分析}

表\ref{tab:findver}报告FinDVer testmini默认阅读器结果。DeepSeek Flash上规划器加回退取得0.701准确率，平均8.76段；固定16段0.660，固定8段0.623。GLM 5.3 Flash上该分支取得0.753，平均9.19段；固定16段0.743，固定8段0.680。Qwen3-8B上该分支取得0.649，平均6.13段；固定16段0.636，固定8段0.629。分支在三个基座上均高于固定16段阅读器，平均少用6.9至9.9段。官方思维链阅读器在DeepSeek Flash和GLM 5.3 Flash上分别取得0.817和0.819（金证据），十段检索版本均为0.787；这些使用不同阅读器和证据接口，是上界参考，不是同条件基线。FinDVer是迁移集，分支未在其上调参。

义务检索分支的主要错误是证据缺失和算术错误。DocFinQA的段落块可能把表格行拆开，导致阅读器找不到分子或分母；规划器的\texttt{missing}字段可以部分修复，但当缺失项不在BM25候选的前32条中时无法恢复。Qwen3-8B在DocFinQA上的精确匹配接近地板，主要原因是回答包含解释性文字，而不是数值本身错误。这些错误指向两个改进方向：更细的表格块边界和更严格的短答案格式约束。

AnchorEvolve-LLM的主要误差来自类别偏置。DeepSeek Flash较多地预测高于区间，纯LLM后验的平衡准确率明显偏低；GLM 5.3 Flash与Qwen3-8B的宏F1更高，但概率质量仍不足以单独承担低损失决策。融合后验在三个基座的准确率、平衡准确率、宏F1和决策损失上均为最高，说明LLM锚点信号与结构化特征并不重复。硬锚点触发较少，当前收益主要来自软锚点的重新加权；全动作读出把多数断言送到二元裁决集，平均集合大小约为2.1，说明后验仍偏保守。共享校准随机森林是强诊断基线，其总体损失与FinPlanRAG接近，但在指导区间内样本上显著更差。

路由器的错误集中在语义边界。DeepSeek Flash和GLM 5.3 Flash把“未来支付金额”“预期股利”等确定问题判为前瞻性断言；这些问题的答案在文件中已经给出，但文本包含未来时态。Qwen3-8B的错误更多，说明小模型的语义边界判断不稳定。对于部署，路由错误比分支错误更容易诊断，因为路由输出带有置信度和理由；系统可以在低置信度时请求人工复核，或回退到保守契约。本文不进一步设计这一回退机制，但结果为后续工作提供了接口。

\begin{table}[H]
\centering
\zihao{5}
\setlength{\tabcolsep}{4pt}
\caption{FinDVer testmini迁移结果（默认阅读器）。固定4段、固定8段和规划器单独运行的结果见正文。}
\label{tab:findver}
\begin{tabular}{llrrr}
\toprule
基座 & 条件 & 准确率 & 宏F1 & 平均段落 \\
\midrule
DeepSeek Flash & 固定16段 & .660 & .735 & 16.00 \\
 & 规划器+回退 & .701 & .759 & 8.76 \\
GLM 5.3 Flash & 固定16段 & .743 & .792 & 16.00 \\
 & 规划器+回退 & .753 & .789 & 9.19 \\
Qwen3-8B & 固定16段 & .636 & .673 & 16.00 \\
 & 规划器+回退 & .649 & .663 & 6.13 \\
\bottomrule
\end{tabular}
\end{table}

% ==================== 第7章 ====================
\chapter{讨论与结论}
\label{ch:conclusion}

\section{讨论}

路由实验回答第一个研究问题。DeepSeek Flash和GLM 5.3 Flash分别达到0.995和0.990的路由准确率，混合效用为0.403和0.419，后者恢复oracle的99.3\%。错路由主要发生在包含未来时态、但答案已经在文件中给出的确定问题。Qwen3-8B的路由准确率降至0.867，混合效用0.263，低于Always-Anchor的0.298。这个对比说明契约识别不是简单的时态匹配：较强基座能够利用查询与截止时点的关系，较小基座则容易把语义边界压缩为表面模式。路由实验因此支持H1在较强基座上成立，同时给出基座能力的边界。

DocFinQA和FinDVer回答第二个研究问题。DeepSeek Flash上，义务检索分支取得0.457精确匹配、平均10.69段，固定16段阅读器取得0.390、平均16段；GLM 5.3 Flash上分别为0.497对0.480，平均10.31段对16段。FinDVer上三个基座分别为0.701、0.753和0.649，均高于固定16段的0.660、0.743和0.636。四个具名RAG适配在DocFinQA上均低于义务检索分支，分支的答案覆盖率0.770和0.810也高于固定预算。规划器单独运行的成本更低，但在GLM 5.3 Flash上留下小缺口，置信度回退在24\%的问题上触发。实验支持H2：类型化义务与覆盖驱动取证在长文档上以更低证据成本取得更高准确率，并可迁移到FinDVer；收益来自证据策略，而不是阅读器或评分口径。

收入指引实验回答第三个研究问题，也给出最需要谨慎解释的结果。纯LLM锚点后验在三个基座上的准确率为0.445、0.454和0.427，决策损失为4.035、3.903和3.586，说明锚点包含非随机信号，但类别偏置和概率校准不足以直接承担低损失决策。融合后验把准确率提高到0.524、0.537和0.537，决策损失降到3.604、3.648和3.577，低于六个具名基线；在GLM 5.3 Flash和Qwen3-8B上，准确率也高于最强具名基线。纯RF诊断取得0.533、0.491、0.486和3.454，是当前最强的数值决策头。这个结果不支持“LLM锚点后验优于强RF”，但支持“语义状态与校准读出分工”：LLM负责路由、阅读和锚点状态，RF负责数值风险控制。H3应表述为融合读出相对具名基线降低决策损失，而不是LLM单独胜过RF。

消融实验进一步说明锚点层的作用与边界。去掉全部锚点后，LLM后验的宏F1从0.233、0.286、0.339降到0.213、0.211、0.272；取消硬锚点掩码后指标几乎不变，说明硬锚点在收入指引上很少触发；逐类去除锚点没有一致赢家，说明锚点类型的作用依赖基座和数据分布。输出形式对照显示，强制点答案的平均损失为4.693，固定三桶全集为4.000，损失感知门为3.454；风险收益来自集合与弃权，而不是更强的点分类器。误差分析显示，DocFinQA的错误主要来自表格块边界与算术，Qwen3-8B的错误主要来自解释性输出，AnchorEvolve的错误主要来自类别偏置，路由错误集中在语义边界。这些结果把改进方向限定在证据表示、锚点过滤和路由信号上，而不是盲目增加模型规模。

实验比较的是同基座适配，不是原系统的未修改端到端运行；标签来自XBRL核验而非人工金标准；混合负载由DocFinQA与收入指引样本构成。这些选择保证了阅读器、证据接口和评分口径可比，也意味着结论适用于给定的数据、基座和评价口径。它们不改变三个研究问题的检验结果，但界定了结果的外推范围。


\section{局限与未来研究}

本文的证据边界集中在三点。第一，基准覆盖仍有限。收入指引主基准只覆盖美国公司的季度营业收入指引；A股回测覆盖40家公司和三个财年，但只到年报、营收、利润、毛利率和披露日期；产能、利用率、良率和排产等更细的数值桥只有模板和前瞻性案例，没有历史端到端基准。混合负载由DocFinQA和收入指引样本构造，真实用户查询的契约边界可能更复杂。

第二，标签来自XBRL核验或披露文件回退，不是人工金标准；A股数值来自公开财务报表，口径差异和重述仍会带来噪声。Qwen3-8B在路由和数值抽取上明显弱于另外两个基座，GLM的数值桥没有触发，说明框架效果依赖基座的契约识别和字段抽取能力。

第三，类型库扩展只在DeepSeek上完成一轮，GLM和Qwen没有候选通过验证门控；全动作读出在总体损失上稳健，但融合与随机森林的总体差异仍不显著。后续工作将补产能与排产的历史基准、更难的路由边界样本和人工审计，并把类型库扩展扩展到多轮和跨基座迁移。

\section{研究结论}

本文提出FinPlanRAG，把金融混合查询的回答形式选择表述为查询契约路由，并以义务检索和AnchorEvolve分别处理确定信息与前瞻性断言。路由器在检索与生成之前判断事实是否已在截止时点存在；义务检索围绕类型化义务形成闭合答案；AnchorEvolve以L0稳定关系、L1行业模式和L2实例锚点维护可更新、可追溯的前瞻性证据状态。校准读出在此基础上生成点答案、可信集或弃权，用于控制最终决策风险。

第一项贡献是查询契约路由。现有问答、断言验证和选择性预测分别固定输出接口，本文把“事实是否已在截止时点存在”置于检索与生成之前，用一次路由决定闭合答案或可校准判断。路由实验支持这一贡献：DeepSeek Flash和GLM 5.3 Flash分别达到0.995和0.990的路由准确率，混合效用接近oracle，并显著高于Always-RAG和Always-Anchor；Qwen3-8B的结果则说明路由能力依赖基座。契约路由的价值不只是提高一个分类指标，而是把混合负载的分析单位从数据集任务改为查询契约，使路由错误和分支错误可以分开诊断。

第二项贡献是义务检索与AnchorEvolve的两条证据分支。义务检索把确定问题绑定为类型化义务，在时点语料上按覆盖取证，由阅读器返回带证据编号的闭合答案；AnchorEvolve把前瞻性断言映射到有限结果空间，用三层锚点、五类资格判定和时点信念更新保留证据的作用域、可靠性和时效。DocFinQA和FinDVer实验支持义务检索的效率主张；收入指引实验和锚点消融支持AnchorEvolve的语义主张。特别是去除全部锚点后LLM后验宏F1下降，说明锚点不是装饰，而是改变了模型判断的中间状态。

第三项贡献是语义状态与校准读出的分工，以及统一的多基准评估。本文在三个数据集、三套基座和六个具名基线上同时检验路由、阅读、锚点推断、后验融合、校准选择和外部迁移，并把分类、校准、证据成本、决策损失和路由效用放在同一口径下报告。需要明确，纯RF是当前最强的数值决策头，融合后验在部分基座上取得更高普通准确率，但决策损失仍高于纯RF；因此本文不主张LLM锚点后验单独胜过强RF。本文主张的是：LLM在路由、阅读和语义状态构造上不可替代，校准读出负责把这些状态转换为低风险决策，二者应当分开评估。

本文的原创性贡献可以概括为三点。第一，把输出契约从固定接口变成查询层的路由变量，使确定问答与前瞻性裁决在同一时点约束下被统一处理。第二，把前瞻性证据从单一标签变成可审计的状态，使证据来源、作用域、资格理由和信念变化可以被逐条检查。第三，把语义判断与风险校准分开，使LLM的语义能力、统计模型的风险控制能力和路由器的契约识别能力可以分别检验、替换和组合。这三点不依赖特定基座或特定行业，构成了本文相对于现有金融RAG和断言验证工作的主要区别。

从金融实践看，框架把答案、证据和风险放在同一条审计链上。分析师对确定问题获得带证据编号的闭合答案，对前瞻性断言获得点判断、可信集或弃权，并可以查看锚点来源和资格理由。机构可以按自身风险偏好替换损失参数和校准后端，而不改变路由与证据状态的结构。生产里程碑、产能投放、监管审批和合同履约等任务与收入指引具有相似的时点结构，本文结果为这些场景提供可复用的接口。

本文的结论有明确边界。收入指引实验采用三桶结果空间，通用的时机与阶段状态仍需在更多前瞻性事件上评估；混合负载由DocFinQA与收入指引样本构成，真实用户查询的语义边界可能更复杂；标签来自XBRL核验而非人工金标准。后续工作将围绕更多前瞻性事件、更稳定的边界路由、更窄的可信集和更严格的人工复核展开。由于路由、锚点抽取和输出读出彼此分离，这些改进可以独立评估，而不需要改变FinPlanRAG的整体接口。


\clearpage
\phantomsection
\addcontentsline{toc}{chapter}{参考文献}
{\zihao{5}\setlength{\baselineskip}{16pt}\begin{thebibliography}{33}
\setlength{\bibsep}{0pt}
\setlength{\itemsep}{0pt}
\setlength{\parsep}{0pt}
\setlength{\bibhang}{2em}

\bibitem[Angelopoulos and Bates(2021)]{angelopoulos2021gentle}
Angelopoulos, A. N., and Bates, S., 2021, ``A Gentle Introduction to Conformal Prediction and Distribution-Free Uncertainty Quantification'', arXiv:2107.07511.

\bibitem[Asai et al.(2024)]{asai2024selfrag}
Asai, A., Wu, Z., Wang, Y., Sil, A., and Hajishirzi, H., 2024, ``Self-RAG: Learning to Retrieve, Generate, and Critique through Self-Reflection'', in \emph{International Conference on Learning Representations}.

\bibitem[Brodersen et al.(2010)]{brodersen2010balanced}
Brodersen, K. H., Ong, C. S., Stephan, K. E., and Buhmann, J. M., 2010, ``The Balanced Accuracy and Its Posterior Distribution'', in \emph{Proceedings of the 20th International Conference on Pattern Recognition}, pp. 3121--3124.

\bibitem[Chen et al.(2026)]{chen2026hcrag}
Chen, S., Tan, H., Li, Y., and Liang, J., 2026, ``HC-RAG: Evidence-Centric Retrieval-Augmented Generation over Heterogeneous Financial Filings'', arXiv:2608.12335.

\bibitem[Chen et al.(2021)]{chen2021finqa}
Chen, Z., Chen, W., Smiley, C., Shah, S., Borova, I., Langdon, D., Moussa, R., Beane, M., Huang, T.-H., Routledge, B., and Wang, W. Y., 2021, ``FinQA: A Dataset of Numerical Reasoning over Financial Data'', in \emph{Proceedings of the 2021 Conference on Empirical Methods in Natural Language Processing}, pp. 3697--3711.

\bibitem[Chi et al.(2026)]{chi2026finsagent}
Chi, J., Tai, Z., Wu, H., Kwok, T. S. T., He, H., Liao, Z., Xiao, B., Jiang, C., Lei, J., Huang, J., Lu, P., Li, M., Ma, L., Wu, Y., Lyu, S., Tian, J., Li, Y., Ma, Y., Guan, S., Hu, D., Cui, Y., Zhou, L., Ding, L., and Wang, X., 2026, ``FinSAgent: Corpus-Aligned Multi-Agent RAG Framework for Evidence-Grounded SEC Filing Question Answering'', arXiv:2607.18102.

\bibitem[Chow(1970)]{chow1970optimum}
Chow, C. K., 1970, ``On Optimum Recognition Error and Reject Tradeoff'', \emph{IEEE Transactions on Information Theory}, Vol. 16, No. 1, pp. 41--46.

\bibitem[Elkan(2001)]{elkan2001cost}
Elkan, C., 2001, ``The Foundations of Cost-Sensitive Learning'', in \emph{Proceedings of the 17th International Joint Conference on Artificial Intelligence}, pp. 973--978.

\bibitem[Geifman and El-Yaniv(2017)]{geifman2017selective}
Geifman, Y., and El-Yaniv, R., 2017, ``Selective Classification for Deep Neural Networks'', in \emph{Advances in Neural Information Processing Systems}, Vol. 30.

\bibitem[Gibbs and Cand\`es(2021)]{gibbs2021adaptive}
Gibbs, I., and Cand\`es, E. J., 2021, ``Adaptive Conformal Inference Under Distribution Shift'', in \emph{Advances in Neural Information Processing Systems}, Vol. 34, pp. 1660--1672.

\bibitem[Guo et al.(2017)]{guo2017calibration}
Guo, C., Pleiss, G., Sun, Y., and Weinberger, K. Q., 2017, ``On Calibration of Modern Neural Networks'', in \emph{Proceedings of the 34th International Conference on Machine Learning}, Vol. 70, pp. 1321--1330.

\bibitem[Hirst et al.(2008)]{hirst2008management}
Hirst, D. E., Koonce, L., and Venkataraman, S., 2008, ``Management Earnings Forecasts: A Review and Framework'', \emph{Accounting Horizons}, Vol. 22, No. 3, pp. 315--338.

\bibitem[Islam et al.(2023)]{islam2023financebench}
Islam, P., Kannappan, A., Kiela, D., Qian, R., Scherrer, N., and Vidgen, B., 2023, ``FinanceBench: A New Benchmark for Financial Question Answering'', arXiv:2311.11944.

\bibitem[Jeong et al.(2024)]{jeong2024adaptiverag}
Jeong, S., Baek, J., Cho, S., Hwang, S. J., and Park, J. C., 2024, ``Adaptive-RAG: Learning to Adapt Retrieval-Augmented Large Language Models through Question Complexity'', in \emph{Proceedings of the 2024 Conference of the North American Chapter of the Association for Computational Linguistics}, pp. 7036--7050.

\bibitem[Jiang et al.(2026)]{jiang2026finrate}
Jiang, Y., Chen, J., Makri, E., Chen, J., Li, P., Maatouk, A., Tassiulas, L., Brenner, E., Xiang, B., and Ying, R., 2026, ``Fin-RATE: A Real-world Financial Analytics and Tracking Evaluation Benchmark for LLMs on SEC Filings'', in \emph{Proceedings of the 32nd ACM SIGKDD Conference on Knowledge Discovery and Data Mining}, pp. 9183--9194.

\bibitem[Jiang et al.(2023)]{jiang2023flare}
Jiang, Z., Xu, F., Gao, L., Sun, Z., Liu, Q., Dwivedi-Yu, J., Yang, Y., Callan, J., and Neubig, G., 2023, ``Active Retrieval Augmented Generation'', in \emph{Proceedings of the 2023 Conference on Empirical Methods in Natural Language Processing}, pp. 7969--7992.

\bibitem[Lewis et al.(2020)]{lewis2020rag}
Lewis, P., Perez, E., Piktus, A., Petroni, F., Karpukhin, V., Goyal, N., Kuttler, H., Lewis, M., Yih, W.-T., Rockt\"aschel, T., Riedel, S., and Kiela, D., 2020, ``Retrieval-Augmented Generation for Knowledge-Intensive NLP Tasks'', in \emph{Advances in Neural Information Processing Systems}, Vol. 33, pp. 9459--9474.

\bibitem[L\'opez de Prado(2018)]{lopezdeprado2018advances}
L\'opez de Prado, M., 2018, \emph{Advances in Financial Machine Learning}, Wiley.

\bibitem[Manski(2003)]{manski2003partial}
Manski, C. F., 2003, \emph{Partial Identification of Probability Distributions}, Springer.

\bibitem[Mansouri et al.(2026)]{mansouri2026finrank}
Mansouri, S., Saad, D., Wahrenburg, M., Weissel, M., and Woebbeking, F., 2026, ``FinRank: An Evidence-Grounded Benchmark for Financial Question Answering and Retrieval over SEC Filings'', arXiv:2608.07400.

\bibitem[Qwen Team(2025)]{qwen2025qwen3}
Qwen Team, 2025, ``Qwen3 Technical Report'', arXiv:2505.09388.

\bibitem[Rangapur et al.(2025)]{rangapur2025finfact}
Rangapur, A., Wang, H., Jian, L., and Shu, K., 2025, ``Fin-Fact: A Benchmark Dataset for Multimodal Financial Fact-Checking and Explanation Generation'', in \emph{Companion Proceedings of the ACM Web Conference 2025}, pp. 785--788.

\bibitem[Reddy et al.(2024)]{reddy2024docfinqa}
Reddy, V., Koncel-Kedziorski, R., Lai, V. D., Krumdick, M., Lovering, C., and Tanner, C., 2024, ``DocFinQA: A Long-Context Financial Reasoning Dataset'', in \emph{Proceedings of the 62nd Annual Meeting of the Association for Computational Linguistics (Volume 2: Short Papers)}, pp. 445--458.

\bibitem[Romano et al.(2020)]{romano2020classification}
Romano, Y., Sesia, M., and Cand\`es, E. J., 2020, ``Classification with Valid and Adaptive Coverage'', in \emph{Advances in Neural Information Processing Systems}, Vol. 33, pp. 3581--3591.

\bibitem[Shinn et al.(2023)]{shinn2023reflexion}
Shinn, N., Cassano, F., Berman, A., Gopinath, A., Narasimhan, K., and Yao, S., 2023, ``Reflexion: Language Agents with Verbal Reinforcement Learning'', in \emph{Advances in Neural Information Processing Systems}, Vol. 36.

\bibitem[Sokolova and Lapalme(2009)]{sokolova2009performance}
Sokolova, M., and Lapalme, G., 2009, ``A Systematic Analysis of Performance Measures for Classification Tasks'', \emph{Information Processing and Management}, Vol. 45, No. 4, pp. 427--437.

\bibitem[Vovk et al.(2005)]{vovk2005algorithmic}
Vovk, V., Gammerman, A., and Shafer, G., 2005, \emph{Algorithmic Learning in a Random World}, Springer.

\bibitem[Xie et al.(2024)]{xie2024finben}
Xie, Q., Han, W., Chen, Z., Xiang, R., Zhang, X., He, Y., Xiao, M., Li, D., Dai, Y., Feng, D., Xu, Y., Kang, H., Kuang, Z., Yuan, C., Yang, K., Luo, Z., Zhang, T., Liu, Z., Xiong, G., Deng, Z., Jiang, Y., Yao, Z., Li, H., Yu, Y., Hu, G., Huang, J., Liu, X.-Y., Lopez-Lira, A., Wang, B., Lai, Y., Wang, H., Peng, M., Ananiadou, S., and Huang, J., 2024, ``FinBen: A Holistic Financial Benchmark for Large Language Models'', in \emph{Advances in Neural Information Processing Systems}, Vol. 37.

\bibitem[Yan et al.(2024)]{yan2024crag}
Yan, S.-Q., Gu, J.-C., Zhu, Y., and Ling, Z.-H., 2024, ``Corrective Retrieval Augmented Generation'', arXiv:2401.15884.

\bibitem[Yao et al.(2023)]{yao2023react}
Yao, S., Zhao, J., Yu, D., Du, N., Shafran, I., Narasimhan, K., and Cao, Y., 2023, ``ReAct: Synergizing Reasoning and Acting in Language Models'', in \emph{International Conference on Learning Representations}.

\bibitem[Zhao et al.(2024)]{zhao2024findver}
Zhao, Y., Long, Y., Jiang, T., Wang, C., Chen, W., Liu, H., Tang, X., Zhang, Y., Zhao, C., and Cohan, A., 2024, ``FinDVer: Explainable Claim Verification over Long and Hybrid-Content Financial Documents'', in \emph{Proceedings of the 2024 Conference on Empirical Methods in Natural Language Processing}, pp. 14739--14752.

\bibitem[Zhou et al.(2026)]{zhou2026fincards}
Zhou, Y., Zhang, F., Chen, Y., Zhang, H., Nakov, P., and Xie, Z., 2026, ``FinCARDS: Card-Based Analyst Reranking for Financial Document Question Answering'', in \emph{Findings of the Association for Computational Linguistics: ACL 2026}, pp. 24836--24852.

\bibitem[Zhu et al.(2021)]{zhu2021tatqa}
Zhu, F., Lei, W., Huang, Y., Wang, C., Zhang, S., Lv, J., Feng, F., and Chua, T.-S., 2021, ``TAT-QA: A Question Answering Benchmark on a Hybrid of Tabular and Textual Content in Finance'', in \emph{Proceedings of the 59th Annual Meeting of the Association for Computational Linguistics}, pp. 3277--3287.

\end{thebibliography}
}

% ==================== 附录 ====================
\appendix
\chapter{完整指标表}
\label{app:fullmetrics}

表\ref{tab:vague-ds}至表\ref{tab:vague-qwen}报告三个基座上的完整十一项指标。主文表\ref{tab:vague}和表\ref{tab:vague-qwen-main}并列报告AnchorEvolve-LLM、包含LLM信号的AnchorEvolve-Fusion以及共享校准RF诊断，以区分语义状态、融合后验与非LLM风险控制读出。完整指标中的AUROC和AUPRC衡量排序质量，ECE和Brier衡量概率校准，Sel@50和AURC衡量选择性预测，决策损失衡量非对称风险。

\begin{table}[htbp]
\centering
\zihao{5}
\setlength{\tabcolsep}{2pt}
\caption{DeepSeek基座上的模糊断言结果（227条XBRL核验断言）。Loss越低越好；Sel@50为50\%覆盖率下的选择性准确率，AURC越低越好。}
\label{tab:vague-ds}
\resizebox{\linewidth}{!}{%
\begin{tabular}{lrrrrrrrrrrr}
\toprule
方法 & Acc & Bal & Macro & MCC & AUROC & AUPRC & ECE & Brier & Loss & Sel@50 & AURC \\
\midrule
FinDVer & \textbf{.551} & .461 & .464 & \textbf{.233} & .674 & .494 & .092 & .580 & 4.678 & .675 & .337 \\
Fin-Fact & .537 & .446 & .450 & .201 & .650 & .475 & .075 & .590 & 4.705 & .675 & .338 \\
Self-RAG & .546 & .451 & .456 & .216 & .670 & .491 & .102 & .580 & 4.564 & \textbf{.684} & \textbf{.333} \\
CRAG & .529 & .423 & .418 & .177 & .643 & .473 & .097 & .593 & 4.802 & .658 & .349 \\
FLARE & .511 & .417 & .411 & .164 & .634 & .467 & .130 & .610 & 4.881 & .596 & .362 \\
Adaptive-RAG & .493 & .406 & .400 & .143 & .584 & .425 & .147 & .634 & 5.093 & .579 & .415 \\
AnchorEvolve-LLM & .445 & .307 & .242 & -0.075 & .465 & .312 & .129 & .666 & 4.035 & .509 & .529 \\
FinPlanRAG（AnchorEvolve-Fusion, $w=0.1$） & .524 & .471 & .468 & .201 & .703 & .503 & \textbf{.033} & \textbf{.564} & 3.604 & .632 & .383 \\
共享校准RF（非LLM诊断） & .533 & \textbf{.491} & \textbf{.486} & .229 & \textbf{.706} & \textbf{.505} & .033 & .566 & \textbf{3.454} & .649 & .383 \\
\bottomrule
\end{tabular}
}
\end{table}


\begin{table}[htbp]
\centering
\zihao{5}
\setlength{\tabcolsep}{2pt}
\caption{GLM 5.3基座上的模糊断言结果（227条XBRL核验断言）。Loss越低越好；Sel@50为50\%覆盖率下的选择性准确率，AURC越低越好。}
\label{tab:vague-glm}
\resizebox{\linewidth}{!}{%
\begin{tabular}{lrrrrrrrrrrr}
\toprule
方法 & Acc & Bal & Macro & MCC & AUROC & AUPRC & ECE & Brier & Loss & Sel@50 & AURC \\
\midrule
FinDVer & .498 & .439 & .442 & .178 & .642 & .439 & .103 & .627 & 5.233 & .605 & .418 \\
Fin-Fact & .454 & .398 & .402 & .089 & .603 & .430 & .178 & .670 & 5.339 & .579 & .427 \\
Self-RAG & .454 & .406 & .404 & .113 & .617 & .434 & .164 & .637 & 5.163 & .588 & .411 \\
CRAG & .515 & .431 & .435 & .173 & .607 & .418 & .148 & .657 & 4.846 & .570 & .441 \\
FLARE & .471 & .388 & .387 & .110 & .583 & .389 & .192 & .705 & 5.286 & .526 & .518 \\
Adaptive-RAG & .427 & .372 & .376 & .059 & .580 & .382 & .227 & .738 & 5.630 & .482 & .542 \\
AnchorEvolve-LLM & .454 & .329 & .295 & -0.004 & .498 & .339 & .145 & .671 & 3.903 & .430 & .552 \\
FinPlanRAG（AnchorEvolve-Fusion, $w=0.1$） & \textbf{.537} & .489 & \textbf{.487} & .225 & .704 & .504 & \textbf{.032} & \textbf{.562} & 3.648 & .640 & \textbf{.376} \\
共享校准RF（非LLM诊断） & .533 & \textbf{.491} & .486 & \textbf{.229} & \textbf{.706} & \textbf{.505} & .033 & .566 & \textbf{3.454} & \textbf{.649} & .383 \\
\bottomrule
\end{tabular}
}
\end{table}


\begin{table}[htbp]
\centering
\zihao{5}
\setlength{\tabcolsep}{2pt}
\caption{Qwen3-8B基座上的模糊断言结果（227条XBRL核验断言）。Loss越低越好；Sel@50为50\%覆盖率下的选择性准确率，AURC越低越好。}
\label{tab:vague-qwen}
\resizebox{\linewidth}{!}{%
\begin{tabular}{lrrrrrrrrrrr}
\toprule
方法 & Acc & Bal & Macro & MCC & AUROC & AUPRC & ECE & Brier & Loss & Sel@50 & AURC \\
\midrule
FinDVer-style & .445 & .422 & .406 & .133 & .564 & .375 & .555 & 1.110 & 5.524 & .456 & .535 \\
Fin-Fact-style & .511 & .477 & .464 & \textbf{.237} & .609 & .404 & .489 & .978 & 4.846 & .491 & .467 \\
Self-RAG-style & .489 & .447 & .437 & .196 & .583 & .394 & .489 & .990 & 5.093 & .482 & .506 \\
CRAG-style & .454 & .418 & .399 & .152 & .566 & .371 & .547 & 1.093 & 5.366 & .456 & .534 \\
FLARE-style & .326 & .298 & .212 & -0.107 & .475 & .322 & .672 & 1.340 & 6.740 & .316 & .664 \\
Adaptive-RAG-style & .344 & .324 & .290 & -0.027 & .492 & .330 & .654 & 1.308 & 6.485 & .298 & .663 \\
AnchorEvolve-LLM & .427 & .353 & .355 & .027 & .562 & .399 & .163 & .647 & 3.586 & .561 & .435 \\
FinPlanRAG（AnchorEvolve-Fusion, $w=0.1$） & \textbf{.537} & .486 & .485 & .227 & \textbf{.707} & \textbf{.509} & \textbf{.027} & \textbf{.562} & 3.577 & .632 & \textbf{.373} \\
共享校准RF（非LLM诊断） & .533 & \textbf{.491} & \textbf{.486} & .229 & .706 & .505 & .033 & .566 & \textbf{3.454} & \textbf{.649} & .383 \\
\bottomrule
\end{tabular}
}
\end{table}


\chapter{Qwen3-8B的DocFinQA结果}
\label{app:qwen}

Qwen3-8B在DocFinQA上返回解释性文字而非短答案，精确匹配接近地板。表\ref{tab:qwen-doc}报告该基座的全部条件；它不用于义务检索分支的主比较，只作为基座能力边界。该结果说明，短答案抽取任务对阅读器的格式遵循能力有额外要求；当阅读器返回解释性文字时，精确匹配和词级F1会同时下降，但答案覆盖率仍可能较高。

\begin{table}[htbp]
\centering
\zihao{5}
\setlength{\tabcolsep}{4pt}
\caption{Qwen3-8B的DocFinQA结果（300题固定样本）。}
\label{tab:qwen-doc}
\begin{tabular}{lrrrr}
\toprule
条件 & 准确率 & 宏F1 & 平均段落 & 覆盖率 \\
\midrule
固定4段 & .100 & .048 & 4.00 & .760 \\
固定8段 & .110 & .040 & 8.00 & .833 \\
固定16段 & .100 & .046 & 16.00 & .853 \\
规划器B=8 & .093 & .034 & 3.67 & .883 \\
规划器+回退 & .093 & .034 & 5.85 & .883 \\
\bottomrule
\end{tabular}
\end{table}

\chapter{命题与完整证明}
\label{app:proofs}

本章给出正文使用的四个命题及完整证明。证明只使用有限结果空间、有限动作集和线性期望，不依赖具体数据或模型。

\section{裁决门的贝叶斯最优性}

\begin{proposition}
给定信念$b\in\Delta(\mathcal H)$、动作集合$\mathcal A=\{\text{answer},\text{adjudicate},\text{abstain}\}$和损失函数$L:\mathcal A\times\mathcal H\to\mathbb R$，选择
\[
\pi^\star=\arg\min_{\pi\in\mathcal A}\mathbb E_{y\sim b}[L(\pi,y)]
\]
是贝叶斯最优动作。
\end{proposition}

\begin{proof}
对任意动作$\pi$，其在信念$b$下的期望损失为
\[
R(\pi,b)=\sum_{y\in\mathcal H}b(y)L(\pi,y).
\]
这是$b$的线性函数，也是有限集合$\mathcal A$上的实值函数。由$\pi^\star$的定义，
\[
R(\pi^\star,b)\le R(\pi,b),\qquad \forall \pi\in\mathcal A.
\]
因此不存在在$b$下期望损失更低的动作。若存在另一个动作$\pi'$满足$R(\pi',b)=R(\pi^\star,b)$，则$\pi'$与$\pi^\star$并列为贝叶斯最优；本文按固定顺序打破平局，不改变最优性。回答动作内部还要在三个结果桶中选择一个标签，但该选择同样是在有限集合上取最小期望损失，因而可以与动作选择合并为同一个贝叶斯决策。
\end{proof}

\section{路由效用的分解}

\begin{proposition}
设确定查询集合为$\mathcal D$，前瞻性查询集合为$\mathcal V$，路由器预测为$\hat{\mathcal D}$和$\hat{\mathcal V}$。混合效用可分解为
\[
U=\frac{1}{n}\sum_{q\in\mathcal D\cap\hat{\mathcal D}}u_D(q)
+\frac{1}{n}\sum_{q\in\mathcal V\cap\hat{\mathcal V}}u_V(q),
\]
其中$u_D(q)$为确定查询的F1，$u_V(q)=1-L(q)/10$为前瞻性查询的效用。
\end{proposition}

\begin{proof}
设$n=|\mathcal D|+|\mathcal V|$。对每条查询定义指示变量
\[
\mathbb I_D(q)=\mathbf 1[q\in\mathcal D],\qquad
\mathbb I_V(q)=\mathbf 1[q\in\mathcal V].
\]
路由器正确时获得对应分支效用，错误时获得0，因此
\[
U=\frac{1}{n}\sum_{q\in\mathcal D}\mathbb I[q\in\hat{\mathcal D}]u_D(q)
+\frac{1}{n}\sum_{q\in\mathcal V}\mathbb I[q\in\hat{\mathcal V}]u_V(q).
\]
把指示变量写成集合交集，即得
\[
U=\frac{1}{n}\sum_{q\in\mathcal D\cap\hat{\mathcal D}}u_D(q)
+\frac{1}{n}\sum_{q\in\mathcal V\cap\hat{\mathcal V}}u_V(q).
\]
oracle路由器对每条查询都正确，故
\[
U_{\text{oracle}}=\frac{1}{n}\sum_{q\in\mathcal D}u_D(q)+\frac{1}{n}\sum_{q\in\mathcal V}u_V(q).
\]
两者之差为错路由项，等于被错路由查询在本分支上的效用之和除以$n$。这一分解说明混合效用可以拆成“路由正确时的分支效用”和“错路由损失”两部分，从而把路由错误与分支错误分开诊断。
\end{proof}

\section{软锚点更新的归一化与单调性}

\begin{proposition}
设$\mathcal H$为有限结果空间，软锚点的隐含桶为$j$，似然为
\[
L_e(k)=
\begin{cases}
p_e, & k=j,\\
(1-p_e)/(|\mathcal H|-1), & k\ne j,
\end{cases}
\]
其中$p_e\in(1/|\mathcal H|,1)$。对任意先验$b(k)>0$和可靠性$\rho\in[0,1]$，更新
\[
b'(k)=\frac{b(k)L_e(k)^\rho}{\sum_{h\in\mathcal H}b(h)L_e(h)^\rho}
\]
是概率分布；当$\rho=0$时$b'=b$，当$\rho=1$时$b'$等于先验与完整似然的归一化乘积；$b'(j)/b'(k)$随$\rho$单调增大。
\end{proposition}

\begin{proof}
分母为正，因为$b(k)>0$且$L_e(k)>0$，所以$b'$非负且和为1，是概率分布。$\rho=0$时$L_e(k)^0=1$，故$b'=b$。对$k\ne j$，
\[
\frac{b'(j)}{b'(k)}
=\frac{b(j)}{b(k)}
\left(\frac{p_e}{(1-p_e)/(|\mathcal H|-1)}\right)^\rho.
\]
括号内的比值大于1（因为$p_e>1/|\mathcal H|$），所以该比值随$\rho$单调增大。可靠性越高，更新越偏向隐含桶；可靠性为0时锚点不改变信念。硬锚点可以看作$\rho\to\infty$并将非可行桶的似然置零的极限，本文用点质量加小底实现，以避免后续软锚点无法恢复。
\end{proof}

\section{可信集的最小基数}

\begin{proposition}
设$b$为$\mathcal H$上的概率分布，$R_\alpha$为满足$\sum_{k\in R}b(k)\ge 1-\alpha$的最小基数子集。把$b$按降序排列为$b_{(1)}\ge\cdots\ge b_{(|\mathcal H|)}$，则$R_\alpha$由前$m$个最大元素组成，其中$m$是使部分和达到$1-\alpha$的最小整数。
\end{proposition}

\begin{proof}
对任意大小为$m$的子集$R$，其概率质量不超过前$m$个最大元素之和。若前$m-1$个元素之和小于$1-\alpha$，则任何大小为$m-1$的子集都无法达到覆盖率，故最小基数至少为$m$。前$m$个元素之和达到$1-\alpha$，因此它们构成一个可行集，且基数等于下界$m$，所以是最小基数可信集。若存在并列的阈值元素，任选其中一个不改变覆盖率；本文按桶的固定顺序打破平局。
\end{proof}

\chapter{LLM提示词}
\label{app:prompts}

本章给出正文三个LLM调用的完整提示词。提示词只使用截止时点之前的输入，不包含已实现结果；JSON字段与正文的结构化输出一致。

\section{路由器提示词}

\begin{lstlisting}[basicstyle=\ttfamily\scriptsize,breaklines=true,columns=fullflexible,frame=single,numbers=none]
You decide whether a financial input is determinate or vague,
then route it to one of two specialists.

DETERMINATE (RAG): the input asks for a fact, number, passage,
estimate, or projection that a filing states or lets you compute
at the cutoff. A question about a future value is determinate when
the filing itself states that value or projection.

VAGUE (ANCHOR): the input asserts or asks about a future outcome
that is not realized at the cutoff. Examples include whether
reported revenue will fall below, within, or above a guidance
range, whether a product will launch, or whether a target will be
met. A management guidance statement is a forward-looking claim,
not a realized fact, even when it contains a number.

Return JSON only:
{"vague":true|false,"route":"rag|anchor",
 "confidence":0.0,"reason":"..."}

Examples:
Input: What was Entergy's net revenue change during 2015?
 -> {"vague":false,"route":"rag"}
Input: Q1 2024 revenue guidance is $12.2-13.2 billion.
 -> {"vague":true,"route":"anchor"}
Input: What did the 10-K report for segment operating profit?
 -> {"vague":false,"route":"rag"}
Input: Will Q3 revenue land within management's guidance range?
 -> {"vague":true,"route":"anchor"}
Input: Management expects second-half revenue of $1.5-1.6 billion.
 -> {"vague":true,"route":"anchor"}

Input: {query}
\end{lstlisting}

\section{锚点抽取提示词}

\begin{lstlisting}[basicstyle=\ttfamily\scriptsize,breaklines=true,columns=fullflexible,frame=single,numbers=none]
You are an AnchorEvolve adjudicator for a vague forward-looking
financial claim. Use ONLY information available at the cutoff.
Do not use the realized outcome.

The claim asks whether reported quarterly revenue will land below,
within, or above the guidance range. The guidance range defines the
three buckets; it is NOT by itself evidence that the outcome will be
within. Use the current-quarter value relative to the guidance
midpoint, the guidance width, and the observed base rates:
- mid_to_current_ratio > 1 means the midpoint requires sequential
  growth; the larger it is, the more likely below.
- mid_to_current_ratio < 1 means the midpoint implies a sequential
  decline; the smaller it is, the more likely above.
- A narrow guidance width makes within more likely; a very wide
  range makes outside buckets more likely.
- Base rates are priors, not destiny; update them with these features.

Extract explicit anchors and return JSON only:
{
  "anchors": [
    {"level":"L2",
     "type":"current_quarter_growth|guidance_width|sector_history|company_history|other",
     "relation":"bound|precedence|lead_time|duration|conservation|rate|unit_conversion|scope",
     "value":"...", "scope":"company|sector|market",
     "reliability":0.0, "hard":false,
     "implied_bucket":"below|within|above|null",
     "rationale":"..."}
  ],
  "bucket_scores":{"below":0.0,"within":0.0,"above":0.0},
  "rationale":"..."
}

Rules: a hard anchor is an explicit constraint that can rule out a
bucket (e.g. a reliable current-quarter value far above the guidance
range); most management guidance is a soft signal. reliability is in
[0,1]. bucket_scores must sum to 1.

Input:
{claim, company, ticker, sector, cutoff, target_year,
 target_quarter, guidance_low_musd, guidance_high_musd,
 guidance_mid_musd, guidance_width_pct, current_q_musd,
 current_q_reliability, observed_bucket_prior, sector_bucket_prior,
 evidence}
\end{lstlisting}

\section{阅读器提示词}

\begin{lstlisting}[basicstyle=\ttfamily\scriptsize,breaklines=true,columns=fullflexible,frame=single,numbers=none]
System:
You are a financial analyst answering a question about a long SEC
filing. Use only the excerpts below. The question may require
arithmetic; do the arithmetic carefully. If the excerpts do not
contain enough information to answer, return "insufficient" and set
sufficient to false.

Return only a JSON object with these keys:
- answer: the final answer as a short string (number or phrase)
- confidence: a number between 0 and 1
- sufficient: true if the excerpts are enough to answer, otherwise false
- missing: a short description of the missing information
- evidence_ids: the paragraph ids you used
- rationale: at most 40 words, no markdown

User:
Use only the excerpts below. Do not use outside knowledge.
This is acquisition round {round}; the excerpts include additional
paragraphs retrieved after you reported missing information.
Question: {question}
Document excerpts:
{evidence}
Return the JSON object now.
\end{lstlisting}

\section{提示词设计说明}

三个提示词共同遵循三条约束。第一，时点约束：路由器只接收查询与截止时点元数据，锚点提取器只接收截止时点之前的证据，阅读器只接收已获取的段落块。第二，结构化输出：三个调用都返回固定字段的JSON，解析失败在评价中计为错误，不进行静默修复。第三，可审计性：锚点提示词要求逐条给出类型、关系、作用域、可靠性、硬软标记、隐含桶和理由；阅读器要求给出证据编号和缺失信息，使答案可以追溯到具体段落或表格块。

\chapter{核心代码}
\label{app:code}

本章给出与正文算法直接对应的核心代码片段，完整实现和运行脚本随论文发布。代码只依赖公开的Python包和本地BM25索引，不包含API密钥。

\section{义务项与覆盖}

\begin{lstlisting}[language=Python]
import re

OBLIGATION_RE = re.compile(
    r"\$?\d[\d,]*(?:\.\d+)?%?|"
    r"\b(?:January|February|March|April|May|June|July|"
    r"August|September|October|November|December)\b",
    flags=re.I,
)

def obligation_terms(text):
    terms, seen = [], set()
    for match in OBLIGATION_RE.findall(str(text or "")):
        term = match.strip().lower().replace("$", "").replace(",", "").replace("%", "")
        if term in {"", "0", "1", "2", "3", "4", "5", "6", "7", "8", "9"}:
            continue
        if term not in seen:
            seen.add(term)
            terms.append(term)
    return terms

def uncovered_terms(terms, chunks):
    text = " ".join(str(getattr(c, "text", "")) for c in chunks)
    lowered = text.lower().replace(",", "")
    covered = set()
    for term in terms:
        if term.isdigit() and len(term) == 4 and term.startswith(("19", "20")):
            if re.search(r"(?<!\d)" + re.escape(term) + r"(?!\d)", lowered):
                covered.add(term)
        elif re.search(r"(?<![0-9.])" + re.escape(term) + r"(?![0-9])", lowered):
            covered.add(term)
    return [t for t in terms if t not in covered]
\end{lstlisting}

\section{锚点信念更新与裁决门}

\begin{lstlisting}[language=Python]
import numpy as np

BUCKETS = ["below", "within", "above"]
HARD_FLOOR = 1e-3
SOFT_P_E = 0.7
LOSS = {"over": 12.0, "under": 8.0, "wrong": 10.0,
        "miss": 8.0, "width": 2.0, "abstain": 5.0}

def normalize(v):
    v = np.maximum(np.asarray(v, dtype=float), 1e-12)
    return v / v.sum()

def belief_from_anchors(prior, parsed):
    b, anchors = prior.copy(), []
    if not parsed:
        return b, anchors
    for anchor in parsed.get("anchors", []) or []:
        implied = str(anchor.get("implied_bucket", "")).lower()
        if implied not in BUCKETS:
            continue
        rel = max(0.0, min(1.0, float(anchor.get("reliability", 0.5) or 0.5)))
        target = BUCKETS.index(implied)
        hard = bool(anchor.get("hard", False)) and rel >= 0.65
        anchors.append({"type": anchor.get("type"),
                        "relation": anchor.get("relation"),
                        "implied_bucket": implied,
                        "reliability": rel, "hard": hard,
                        "scope": anchor.get("scope")})
        if hard:
            b = np.array([1.0 if i == target else HARD_FLOOR for i in range(3)])
        else:
            like = np.full(3, (1.0 - SOFT_P_E) / 2.0)
            like[target] = SOFT_P_E
            b = b * np.power(like, rel)
        b = normalize(b)
    scores = parsed.get("bucket_scores")
    if isinstance(scores, dict):
        marginal = normalize(np.array([float(scores.get(x, 0.0)) for x in BUCKETS]))
        b = normalize(0.5 * b + 0.5 * marginal)
    return b, anchors

def credible_set(b, alpha=0.9):
    order = np.argsort(-b)
    out, total = [], 0.0
    for idx in order:
        out.append(int(idx))
        total += float(b[int(idx)])
        if total >= alpha:
            break
    return sorted(out)

def gate(b, alpha=0.9):
    exp_ans = [sum(float(b[y]) * (0.0 if a == y else
               LOSS["over"] if (a, y) == (2, 0) else
               LOSS["under"] if (a, y) == (0, 2) else
               LOSS["wrong"]) for y in range(3)) for a in range(3)]
    a = int(np.argmin(exp_ans))
    R = credible_set(b, alpha)
    exp_set = sum(float(b[y]) * ((0.0 if y in R else LOSS["miss"]) +
                 LOSS["width"] * max(0, len(R) - 1)) for y in range(3))
    if exp_ans[a] <= exp_set and exp_ans[a] <= LOSS["abstain"]:
        return {"mode": "answer", "a": a, "R": R}
    if exp_set <= LOSS["abstain"]:
        return {"mode": "adjudicate", "a": a, "R": R}
    return {"mode": "abstain", "a": a, "R": R}
\end{lstlisting}

\section{校准随机森林}

\begin{lstlisting}[language=Python]
import numpy as np
from sklearn.ensemble import RandomForestClassifier

for fold in range(FOLDS):
    train_rows = [row for row in rows
                  if fold_of(company_key(row[2]), seed) != fold
                  and row[1] < boundary]
    fold_test = [row for row in test_rows
                 if fold_of(company_key(row[2]), seed) == fold]
    X_train, C_train, y_train, meta_train = build(train_rows, by_company)
    X_test, C_test, y_test, meta_test = build(fold_test, by_company)
    train_features = np.hstack([X_train, C_train])
    test_features = np.hstack([X_test, C_test])
    model = make_rf(seed, class_weight, min_samples_leaf)
    model.fit(train_features, list(map(str, y_train)))
    probabilities = align(model.predict_proba(test_features), model.classes_)
    groups = np.asarray([company_key(record) for record in meta_train])
    oof, truths = [], []
    labels = np.asarray(list(map(str, y_train)))
    for fit_index, cal_index in GroupKFold(n_splits=min(3, len(set(groups)))).split(
        train_features, labels, groups
    ):
        inner = make_rf(seed, class_weight, min_samples_leaf)
        inner.fit(train_features[fit_index], labels[fit_index].tolist())
        oof.append(align(inner.predict_proba(train_features[cal_index]),
                         inner.classes_))
        truths.extend(labels[cal_index].tolist())
    temperature = fit_temperature(np.vstack(oof), truths)
\end{lstlisting}

\chapter{实现结构与复现边界}
\label{app:implementation}

\section{实现结构}

实现分为五个模块。路由模块读取原始查询和截止时点元数据，返回路由标签、置信度和简短理由。锚点模块读取断言、指引区间、当前季度值、基率和证据项，返回L2锚点及其可靠性、硬软标记和隐含桶。义务检索分支模块建立时点段落与表格索引，按义务覆盖迭代取证，并调用阅读器返回闭合答案。校准模块在开发集与验证集上训练随机森林并做温度缩放。汇总模块把逐条预测、路由决策和指标写入统一结果文件。各模块通过统一的数据结构连接，LLM调用按请求负载的哈希缓存，因此相同输入可以复现相同输出。

\section{复现边界}

公开结果包含逐条预测、路由决策和汇总指标；LLM响应缓存不随论文提交。随机森林的折种子、校准方式和融合权重选择在正文中给出。所有API密钥从环境变量读取，不写入结果文件或论文源文件。具名基线为同基座适配，代码、提示词和逐条输出随论文发布；原始系统在其原始数据上的未修改端到端结果不在本文的复现范围内。

